\documentclass[a4paper,12pt] {article}

\usepackage[utf8]{inputenc}
\usepackage[english,russian]{babel}

\usepackage{geometry}
\geometry{left=3cm}
\geometry{right=1cm}
\geometry{top=2cm}
\geometry{bottom=2cm}
\linespread{1.3}

\usepackage{graphicx}
\graphicspath{{pic/}}
\DeclareGraphicsExtensions{.png,.jpg}

\newcommand{\tocsecindent}{\hspace{7mm}}


\begin{document}

\begin{titlepage}
		
\thispagestyle{empty}
		
  \begin{center}
  
   
   ПРАВИТЕЛЬСТВО РОССИЙСКОЙ ФЕДЕРАЦИИ \\
   
   \vspace{0.25cm}
   ФЕДЕРАЛЬНОЕ ГОСУДАРСТВЕННОЕ БЮДЖЕТНОЕ
   ОБРАЗОВАТЕЛЬНОЕ УЧРЕЖДЕНИЕ ВЫСШЕГО ОБРАЗОВАНИЯ 
   «САНКТ-ПЕТЕРБУРГСКИЙ ГОСУДАРСТВЕННЫЙ УНИВЕРСИТЕТ» 
   (СПбГУ)
   
   \vspace{0.25cm}
   Направление: Прикладные математика и физика
   
 
   \begin{figure}[h!]
   	\begin{center}
   		\includegraphics[width=0.25\linewidth]{Gerb}
   	\end{center}
   \end{figure}


     \textbf{\Large Полуавтоматическое структурирование
     изображений в социальной сети с помощью 
     методов машинного обучения}
     
 \bigskip

\vfill

\begin{flushright}
    \footnotesize{
    Выполнил: студент 2-го курса магистратуры кафедры вычислительной физики, А. Р. Шабанов
    
    Руководитель: доцент кафедры вычислительной физики, к.ф.-м.н., А. С. Немнюгин
    
    Рецензент: разработчик интеллектуальных систем в ООО “Мэйл.ру”, А. А. Дзюба
    }
\end{flushright}

\end{center}

\vfill

\newlength{\ML}
\settowidth{\ML}{«\underline{\hspace{0.7cm}}» \underline{\hspace{2cm}}}


\begin{center}
    Санкт - Петербург \\
    2019
\end{center}

\end{titlepage}




\tableofcontents
\setcounter{page}{2}


\newpage
\section*{Введение}
\addcontentsline{toc}{section}{\tocsecindent{Введение}}
\indent
\indent
С каждым днем пользователи социальных сетей создают и потребляют все б\'{о}льшие
и б\'{о}льшие объемы информации, в том числе огромное количество фотографий.
Построение системы для быстрой и точной навигации в миллионах изображений
 --- не тривиальная задача. Одно из самых распространенных решений этой
 проблемы заключается в использовании тегов. Частный случай такого подхода ---
 использование хештегов в социальных сетях.


\indent
\textit{Хештег} --- это любое слово или фраза без пробелов, перед которой стоит 
символ \#, который называется \textit{диез} или \textit{решетка}, а в англоязычном 
варианте -- \textit{hash}, отсюда и название. Приведем несколько примеров:
\htag{masterwork}, \htag{spbu}, \#htag{lovesport}. Обычно в браузере или 
приложении хештег отображается как гипертекст, кликнув по которому можно 
получить список публикаций, снабженных таким же тегом.

\indent
Кроме простоты и удобства теги обладают еще одним полезным свойством 
-- они позволяют не думать об  иерархии структурируемой информации. 
Например, набор изображений можно разложить
по папкам, создав иерархию по датам, геолокациям или авторству. Причем в отдельных
случаях подобрать наиболее подходящую иерархию бывает затруднительно. Проблему
можно решить так: достаточно поставить несколько тегов для всех изображений,
а сами они могут храниться в плоской системе файлов.
Благодаря этому свойству тегирование используются  для рубрикации контента 
не только только онлайн, но и в оффлайн приложениях, например, 
просмоторщиках фотографий. 


\indent
Можно выделить две разновидности популярных хэштегов в социальных сетях.
Первые, используемые недолго и посвященные каким-то социальным явлениям
или событиям, например: \htag{elections2018}, \htag{metoo}. 
И вторые, широкораспространенные, но не связанные с новостной повесткой, 
например: \htag{sport}, \htag{cafe}; они и будут нас интересовать.
 Данная работа посвящена разработке интеллектуальной системы, 
подсказывающей пользователю релевантные хештеги к загружаемым фотографиям.
 Кроме того, с помощью такой системы можно решать и 
  \textquotedblleft обратную\textquotedblright  задачу -- определять,
уместно ли поставлены те или иные теги к заданным изображениям. 
Способность системы давать ответ на такой вопрос можно использовать для 
выявления злоупотреблений со стороны пользователей. Например, зачастую
 в рекламных целях продвигаемые публикации снабжают множеством популярных 
 тегов, не имеющих никакого отношения к публикуемой информации.
 
  



\newpage
\section{Постановка задачи}



\section{Используемые инструменты}



\section{Сверточные нейронные сети}

\subsection{Введение в теорею сверточных нейросетей}

\subsection{Обзор используемых архитектур}



\section{Работа с данными}

\subsection{Датасет SUN}

\subsection{База WordNet}

\subsection{Адоптация датасета SUN}



\section{Численные эксперименты}

\subsection{Архитектура программы}

\subsection{Особенности реализации}

\subsection{Полученные результаты}

\subsection{Обсуждение результатов}



\section{Валидация результатов}



\section*{Выводы}
\addcontentsline{toc}{section}{\tocsecindent{Выводы}}



\section*{Список литературы}
\addcontentsline{toc}{section}{\tocsecindent{Список литературы}}


\end{document} 
